\scanpage{183}
les fibres \unsure{de} $X/S$ soient \unsure{abs.|art.|rel.} \deleted{intègres}\inserted{primaires}. Alors on voit que l'on a

\margin{\uline{Prop}}
\[
\underline{\underline{\unsure{\operatorname{Dix}|\operatorname{Div}}}}_{X/S}^{\underline L}\simeq\underline P(\underline E)
\]
\uline{fibré projectif associé à} $\underline E$. Comme ce dernier est aussi isomorphe à $\operatorname{Proj}(\deleted{\unsure{\bigoplus|\underline E}}{\textstyle\coprod}_n\operatorname{Sym}^{n}(\underline E))$, \unsure{et|est} à fibres simples, on voit ceci :

\uline{Prop} Supposons $X/S$ \uline{primaire}. Alors \deleted{$\underline{\underline{\operatorname{Dix}}}^{L}_{X/S}$ est} les conditions suivantes sont équivalentes :

(i) $\underline{\underline{\operatorname{Div}}}_{X/S}^{L}$ est plat sur $S$ en $s$

% slanted
\margin{\deleted{\unsure{Doit}} \unsure{Dégager dès|Diagramme dès} $\hookrightarrow$ \unsure{que le mm sur} les fibrés projectifs}
(ii) $\underline{\underline{\operatorname{Div}}}_{X/S}^{L}$ est \uuline{simple} sur $S$ en $s$

(iii) Pour $n$ grand, $\operatorname{Sym}^{n}(\underline E)$ est loc libre sur $S$ en $s$

(iv) Pour \inserted{$k$ convenable,} $n$ grand, $\operatorname{Sym}^{nk}(\underline E)$ est loc. libre sur $S$ \inserted{en $s$} pour tt $n$ grand.

\diagram{\begin{tikzcd}[ampersand replacement=\&,row sep=tiny]
{\text{(i)}}\arrow[dd,Rightarrow,bend right=22]\\
{\text{(ii)}}\arrow[u,Rightarrow]\\
{\text{(iii)}}\arrow[d,Rightarrow]\\
{\text{(iv)}}\arrow[uuu,Rightarrow,bend left=44]
\end{tikzcd}}
(v) $\underline E$ est \deleted{libre} loc libre sur $S$ en $s$

\margin{par déf. $\unsure{\Updownarrow}$}
(vi) $\underline L$ est \unsure{co-plat|$0$-plat} rel. $S$

\margin{par Künneth}
(vii) $\underline H_f^0(\underline L)_{\inserted{s}}\deleted{\unsure{\otimes|\otimes k(s)}}\longrightarrow H^0(X_s,\underline L_{(s)})$ est \uuline{surjectif}
% archive sheet number: 182
